%TV text
{\Large {\bfseries Т}еле{\bfseries б}ачення ({\bfseries TV})}
{
\definecolor{HumanAudioColor}{rgb}{0.2,0.2,0.6}
\begin{itemize}

\item Ефірне телебачення використовує діапазони VHF та UHF (30MHz - 3GHz)  \psframebox[linestyle=none]{\rput(0.05,.05){\psframebox[fillstyle=solid,framearc=0.25,fillcolor=red]{\textcolor{white}{\tiny 05}}}}\hspace{.1in}.

\item Супутникове телебачення передається в C-діапазоні (4 - 8 GHz)
\psframebox[linestyle=none]{\rput(0.05,.05){\psframebox[fillstyle=solid,framearc=.25,fillcolor=green]{\textcolor{black}{\tiny 05}}}}\hspace{.1in} та Ku-діапазоні (12 - 18 GHz)  \psframebox[linestyle=none]{\rput(0.04,0.02){\psdots[framesep=1pt,fillcolor=BrightGreen,fillstyle=solid,dotstyle=triangle,linecolor=BrightGreen](0,0)}}\hspace{.08in}, де показано один з багатьох супутників. Для усунення перешкод станції транслюються з чергуванням поляризацій, наприклад, канал 1 --- вертикальна, а канал 2 --- горизонтальна, і навпаки на сусідніх супутниках.

%\item Each station is transmitted in its own band.
\item Телеканали, що передаються через кабельне телебачення (CATV), показані як  \psframebox[linestyle=none]{\rput(0.05,.05){\psframebox[fillstyle=solid,framearc=0.1,fillcolor=blue]{\textcolor{white}{\tiny 05}}}}\hspace{.1in}. Канали CATV, що починаються з ``T-'', --- це канали зворотного зв'язку до кабельної телестанції (наприклад, новинні канали).
\item
\begin{tabular}[t]{@{}l@{\hspace{0.1in}}r}
\begin{minipage}[t]{2.8in}
Ефірні та кабельні аналогові телестанції транслюються з окремими несучими частотами відео, кольору та аудіо, згрупованими разом у канальну смугу наступним чином:
\end{minipage}&
\begin{minipage}[t]{2.4in}
	\rput(0,-.17in){
	\psframebox[linestyle=none]{\psset{xunit=1.25in}%was 1.7
		\psline{|<->|}(0,.28)(2.5,.28)\uput{1pt}[90](1.25,.28){6MHz}
		\white
		\psframe[linestyle=solid,framearc=0.25,fillcolor=red,fillstyle=solid](0,0)(2.5,.25)
		\psdots[linecolor=white,dotstyle=triangle*](0.625,0.02)(1.915,0.02)(2.375,0.02)
		\psline[linecolor=white](0.625,0)(0.625,.25)
		\psline[linecolor=white](1.915,0)(1.915,.15)
		\psline[linecolor=white](2.375,0)(2.375,.25)
		\psline[linecolor=white]{<->}(0,.125)(.625,.125)\rput(.29,.125){
			\psframebox[linearc=1pt,linestyle=none,framesep=0pt,fillcolor=white,fillstyle=solid]{\textcolor{Black}{1.25MHz}}}
		\psline[linecolor=white]{<->}(.625,.07)(1.915,.07)\rput(1.27,.07){
			\psframebox[linearc=1pt,linestyle=none,framesep=0pt,fillcolor=white,fillstyle=solid]{\textcolor{Black}{3.58MHz}}}
		\psline[linecolor=white]{<->}(.625,.18)(2.375,.18)\rput(1.50,.18){
			\psframebox[linearc=1pt,linestyle=none,framesep=0pt,fillcolor=white,fillstyle=solid]{\textcolor{Black}{4.5MHz}}}
		\uput{1pt}[270](.625,0){Відео}
		\uput{1pt}[270](1.915,0){Колір}
		\uput{1pt}[270](2.375,0){Аудіо}
		}
	}
\end{minipage}
\end{tabular}
%TV horizontal refresh
\item Сигнал горизонтальної розгортки 15,7 kHz є поширеним забруднювачем при прослуховуванні VLF \hspace{.05in}
  \psframebox[framearc=0,linearc=0]{\rput(0,.05){
	\psframe[cornersize=relative,linecolor=white, linestyle=solid, linewidth=0.8pt,fillstyle=solid,framearc=.25,fillcolor=red,linearc=0.25](-.1,-.1)(.1,.1)
	\psline[linecolor=white,linestyle=solid,linewidth=1pt]{<->}(-.1,0)(.1,0)
  }}\hspace{0.07in}.
\item Методи цифрового стиснення використовуються для HDTV-трансляцій, щоб вмістити більше каналів у ту ж смугу 6MHz, що й аналогове телебачення.

\end{itemize}

}



